% Chapter 4, Topic _Linear Algebra_ Jim Hefferon
%  http://joshua.smcvt.edu/linearalgebra
%  2001-Jun-12
\topic{Populasi Stabil}
\index{populasi stabil|(}
\index{populasi!stabil|(}
Bayangkan sebuah taman suaka bagi satwa dari suatu spesies yang hendak kita
lindungi. Taman itu tidak berpagar sehingga satwa melintasi batasnya, baik dari
dalam ke luar maupun dari luar ke dalam. Setiap tahun, $10\%$ satwa di dalam
taman pergi dan $1\%$ satwa di luar taman masuk. Dapatkah dicapai tingkat
populasi yang stabil? Adakah populasi di dalam taman dan di seluruh wilayah
lain yang tetap konstan dari waktu ke waktu, dengan jumlah satwa yang pergi
sama dengan jumlah satwa yang masuk?

Misalkan $p_n$ adalah populasi di dalam taman pada tahun ke-$n$ dan $r_n$
adalah populasi di seluruh wilayah lain.
\begin{align*}
  p_{n+1} 
  &=.90p_n+.01r_n    \\
  r_{n+1}
  &=.10p_n+.99r_n 
\end{align*}
Kita memperoleh persamaan matriks berikut.
\begin{equation*}
  \colvec{p_{n+1} \\ r_{n+1}}
  =\begin{mat}[r]
    .90  &.01  \\
    .10  &.99
  \end{mat}
  \colvec{p_{n} \\ r_{n}}
\end{equation*}
Populasi akan stabil jika $p_{n+1}=p_n$ dan $r_{n+1}=r_n$, sehingga persamaan
matriks $\vec{v}_{n+1}=T\vec{v}_{n}$ menjadi $\vec{v}=T\vec{v}$.
Jadi, kita mencari vektor eigen $T$ yang terkait dengan nilai eigen
$\lambda=1$. Persamaan $\zero=(\lambda I-T)\vec{v}=(I-T)\vec{v}$ adalah
\begin{equation*}
  \begin{mat}[r]
      0.10  &-0.01  \\
      -0.10  &0.01
  \end{mat}
  \colvec{p \\ r}=
  \colvec[r]{0 \\ 0}
\end{equation*}
dan memberikan ruang eigen yang terdiri atas vektor-vektor dengan syarat
$p=.1r$. Sebagai contoh, jika mula-mula populasi di dalam taman adalah
$p=10\,000$ satwa dan populasi di seluruh wilayah lain adalah $r=100\,000$
satwa, maka setiap tahun sepuluh persen satwa di dalam taman pergi (seribu
satwa), sedangkan satu persen satwa dari luar masuk (juga seribu satwa).
Populasi tersebut stabil dan mempertahankan dirinya sendiri.

Sekarang bayangkan bahwa kita hendak meningkatkan populasi total spesies ini
sebesar $1\%$ per tahun. Dalam suatu pengertian, tingkat populasinya akan stabil
secara dinamis, berbeda dari tingkat populasi statis pada kasus $\lambda=1$.
Persamaan $\vec{v}_{n+1}=1.01\cdot\vec{v}_n=T\vec{v}_{n}$ mengarah pada
$(1.01 I-T)\vec{v}=\zero$, yang memberikan sistem berikut.
\begin{equation*}
  \begin{mat}[r]
     0.11   &-0.01  \\
     -0.10  &0.02
  \end{mat}
  \colvec{p \\ r}=
  \colvec[r]{0 \\ 0}
\end{equation*}
Matriks ini nonsingular sehingga satu-satunya solusi adalah $p=0$, $r=0$.
Jadi, tidak ada populasi awal tak trivial yang menghasilkan laju pertumbuhan
tahunan tetap sebesar satu persen pada $p$ dan $r$.

Kita dapat mencari laju yang memungkinkan suatu populasi awal menghasilkan
perilaku pertumbuhan yang tetap. Tinjau $\lambda\vec{v}=T\vec{v}$ dan
selesaikan untuk $\lambda$.
\begin{equation*}
  0=\begin{vmat}
    \lambda-.9  &-.01  \\
    -.10        &\lambda-.99
  \end{vmat}
  =(\lambda-.9)(\lambda-.99)-(.10)(.01)
  =\lambda^2-1.89\lambda+.89
\end{equation*}
Kita telah mengetahui bahwa $\lambda=1$ merupakan salah satu solusi persamaan
karakteristik ini. Solusi lainnya adalah $0.89$. Namun, ruang eigen untuk
$\lambda=0.89$ memenuhi $p=-r$, sehingga tidak memuat populasi nonnegatif tak
trivial. Karena itu, satu-satunya keadaan stabil yang bermakna sebagai populasi
adalah kasus $\lambda=1$, dengan populasi total yang tidak bertambah atau
berkurang.
 
Jadi, salah satu cara memandang nilai eigen dan vektor eigen adalah sebagai
pemberi keadaan stabil suatu sistem. Jika nilai eigennya satu, sistem tersebut
statis; jika nilainya bukan satu, kestabilannya bersifat dinamis---asalkan
vektor eigennya memang berada dalam domain keadaan yang bermakna.





\begin{exercises}
  \item 
    Untuk taman yang dibahas di atas, berapakah populasi awal di dalam taman
    jika kedua populasi hendak berkurang sebesar $11\%$ setiap tahun?
    \begin{answer}
      Persamaan 
      \begin{equation*}
         0.89I-T
         =
         \begin{mat}
           0.89 &0 \\
           0 &0.89
         \end{mat}
         -
         \begin{mat}
           .90  &.01  \\
           .10  &.99
         \end{mat}
         =
         \begin{mat}
           -.01  &-.01  \\
           -.10  &-.10
         \end{mat}
      \end{equation*} 
      mengarah pada sistem berikut.
      \begin{equation*}
        \begin{mat}
          -.01  &-.01  \\
          -.10  &-.10
        \end{mat}
        \colvec{p \\ r}
        =
        \colvec{0  \\ 0}
      \end{equation*}
      Jadi, vektor eigennya memenuhi $p=-r$. Tidak ada populasi awal nonnegatif
      tak trivial yang memenuhi syarat ini; satu-satunya keadaan populasi yang
      mungkin adalah $p=r=0$.
    \end{answer}
  \item 
    Apa yang terjadi pada populasi taman jika populasi di luar taman bertambah
    sebesar $1\%$ per tahun?
    Apakah laju pertumbuhan taman tertinggal atau mendahului laju di luar taman?
    Andaikan populasi awal taman sepuluh ribu dan populasi awal di luar taman
    seratus ribu, lalu hitung selama sepuluh tahun.
    \begin{answer}
      Dengan memperlakukan populasi luar sebagai besaran yang diberikan oleh
      $r_n=100\,000(1.01)^n$ dan menggunakan
      $p_{n+1}=0.90p_n+0.01r_n$, \Sage{} memberikan hasil berikut.
\begin{lstlisting}
sage: p, r0 = 10000.0, 100000.0
sage: for y in range(11):
....:     r = r0*(1.01)^y
....:     print(y, p, r)
....:     p = 0.90*p + 0.01*r
0 10000.0 100000.0
1 10000.0 101000.0
2 10010.0 102010.0
3 10029.1 103030.1
4 10056.491 104060.401
5 10091.44591 105101.00501
6 10133.311369 106152.01506
7 10181.500383 107213.535211
8 10235.485697 108285.670563
9 10294.793833 109368.527268
10 10358.999722 110462.212541
\end{lstlisting}
      Jadi, selama sepuluh tahun populasi di dalam taman bertambah sekitar
      $3.59\%$, sedangkan populasi di luar taman bertambah sekitar $10.46\%$.
      Laju pertumbuhan di dalam taman tertinggal.
    \end{answer}
  \item 
    Taman yang dibahas di atas kini dipagari sebagian, sehingga setiap tahun
    hanya $5\%$ satwa di dalam taman yang pergi (sementara sekitar $1\%$ satwa
    dari luar tetap masuk). Dalam keadaan apa taman tersebut sekarang dapat
    mempertahankan populasi yang stabil?
    \begin{answer}
      Kita menginginkan populasi taman tahun berikutnya merupakan gabungan
      satwa yang tetap tinggal dan satwa yang masuk dari luar,
      $p_{n+1}=0.95p_n+0.01r_n$.
      Populasi tahun berikutnya di wilayah lain merupakan gabungan satwa yang
      meninggalkan taman dan satwa yang tetap berada di luar,
      $r_{n+1}=0.05p_n+0.99r_n$. 
      Persamaan matriks 
      \begin{equation*}
        \colvec{p_{n+1}  \\ r_{n+1}}
        =
        \begin{mat}
          0.95  &0.01  \\
          0.05  &0.99
        \end{mat}
        \colvec{p_n  \\ r_n}
      \end{equation*}
      berarti bahwa untuk menentukan nilai eigen kita harus menyelesaikan
      persamaan berikut.
      \begin{equation*}
        0
        =
        \begin{vmat}
          \lambda-0.95  &-0.01  \\
          -0.05         &\lambda-0.99
        \end{vmat}
        =\lambda^2-1.94\lambda+0.94
      \end{equation*}
      \Sage{} memberikan hasil berikut.
\begin{lstlisting}
sage: x = var('x')
sage: qe = (x^2 - 1.94*x + .94 == 0)
sage: print solve(qe, x)
[
x == 47/50,
x == 1
]
\end{lstlisting}
      Untuk nilai eigen $1$, persamaan $(I-T)\vec v=\zero$ memberikan
      $0.05p=0.01r$, atau $r=5p$. Jadi, populasi stabil ketika populasi di luar
      taman lima kali populasi di dalam taman. Nilai eigen $0.94$ kembali
      menghasilkan $p=-r$ dan tidak memberikan populasi nonnegatif tak trivial.
    \end{answer}
  \item 
    Andaikan suatu spesies burung hanya hidup di Kanada, Amerika Serikat, atau
    Meksiko. Setiap tahun, $4\%$ burung Kanada berpindah ke Amerika Serikat dan
    $1\%$ berpindah ke Meksiko. Setiap tahun, $6\%$ burung Amerika Serikat
    berpindah ke Kanada dan $4\%$ ke Meksiko. Dari Meksiko, setiap tahun $10\%$
    berpindah ke Amerika Serikat dan $0\%$ ke Kanada.
    \begin{exparts}
      \partsitem Berikan matriks transisinya.
      \partsitem Adakah keadaan yang membuat populasi di ketiga negara
         tersebut konstan?  
    \end{exparts}
    \begin{answer}
      \begin{exparts}
        \partsitem Berikut rekurensinya.
          \begin{equation*}
            \colvec{c_{n+1} \\ u_{n+1} \\ m_{n+1}}
            =
            \begin{mat}
              .95  &.06  &0  \\
              .04  &.90  &.10  \\
              .01  &.04  &.90
            \end{mat}
            \colvec{c_n \\ u_n  \\ m_n}
          \end{equation*}
        \partsitem Sistem
          \begin{equation*}
            \begin{linsys}{3}
              .05c &- &.06u &  &     &= &0  \\
             -.04c &+ &.10u &- &.10m &= &0  \\
             -.01c &- &.04u &+ &.10m &= &0
            \end{linsys}
          \end{equation*}
          memiliki tak hingga banyak solusi.
\begin{lstlisting}
sage: var('c,u,m')
(c, u, m)
sage: eqns = [ .05*c-.06*u  == 0,
....:         -.04*c+.10*u-.10*m == 0,
....:         -.01*c-.04*u+.10*m == 0 ]
sage: solve(eqns, c,u,m)
[[c == 30/13*r1, u == 25/13*r1, m == r1]]            
\end{lstlisting}
      \end{exparts}
    \end{answer}
\end{exercises}
\index{populasi!stabil|)}
\index{populasi stabil|)}
\endinput



